\begin{afterword}
\summaryhead

The essay lists the habits of mind that the author thinks make up good thinking:
\begin{itemize}
\item curiosity: wanting to know;
\item relinquishment: giving up beliefs the truth destroys;
\item lightness: changing your mind quickly;
\item evenness: holding favored and unfavored ideas to the same standard;
\item argument: discussing honestly with others;
\item empiricism: tying beliefs to observation;
\item simplicity: avoiding needless detail;
\item humility: preparing for your own errors;
\item perfectionism: never settling for almost right;
\item precision: making narrow claims that can be tested;
\item scholarship: learning many fields.
\end{itemize}
A twelfth, unnamed virtue comes before all the others: aiming at the true answer rather than at following rules. The author illustrates it with Musashi's \textsc{Book of Five Rings}, where every movement of the sword must aim to cut. The whole essay is written as a series of commands and maxims.

\responsehead

The essay's strength is its ambition. It tries to say, in a few pages, what intellectual honesty consists of, and several of its sentences are very good. ``To confess your fallibility and then do nothing about it is not humble; it is boasting of your modesty.'' The two guesses, between 1 and 100 and between 40 and 50, make the value of precision concrete. The Great Teacher who says the sky is green makes a sharp point about authority.

Its form is also its weakness. Commands cannot be checked, and the essay rarely argues for anything. It tells readers to tie every belief to observation, then ends by promising that years of practice will let them ``glimpse the center,'' without saying what that would look like. Some of the virtues pull against each other: lightness against precision, humility against perfectionism. A careful essay would say how to balance them, instead of issuing each command at full strength.

Lists of virtues and vices of the mind are old. Aristotle separated intellectual virtues from moral ones in the \textsc{Nicomachean Ethics}. Francis Bacon's \textsc{Novum Organum} (1620) catalogued the ``idols'' that distort human understanding, which makes it a close ancestor of this essay's warnings. Since Ernest Sosa's 1980 paper ``The Raft and the Pyramid,'' philosophers have studied the subject under the name ``virtue epistemology.'' Richard Feynman put much of the list into one sentence in a 1974 commencement address at Caltech: ``The first principle is that you must not fool yourself—and you are the easiest person to fool.''

The essay was first published on the author's website, before LessWrong existed (the site opened in 2009). It reads like a charter for the community that later gathered there.

As writing, it has many memorable lines, too many metaphors, and a solemn tone. A teacher would admire the ambition, circle the contradictions, and ask the student to write one paragraph of argument for every paragraph of commands.
\end{afterword}
